%% ==========================================================================
%%  FINAL REPORT --- SECTION 10, CONCLUSION
%%  The survey's conclusion.tex concludes the survey and is left untouched.
%% ==========================================================================
\section{Conclusion}
\label{sec:conclusion}

\nudgeplain{} trains a recommender privately and stops at the scores.
\pirsonaplain{} delivers privately with a training core the field has moved past.
\sys{} composes the two and measures the result, and three findings stand out.

First, the composition is cheap in the way that matters to its designers and
expensive in a way they did not measure. It costs exactly the sum of its halves,
private training costs 96\,ms per user-session on a wide-area link when amortised
over 943 users, and the normaliser that reveals nothing loses only a few
thousandths of nDCG@20. But the rounds are dominated by truncation, not by the
comparison gate that looks like the hard part, and delivery as built pays a round
trip per recommendation, which lets a full catalogue download beat it on a
wide-area link.

Second, the private read is necessary, and the cheaper alternatives do not
replace it. Because the item embeddings are public by design, ten fetches seen in
the clear recover a user's taste to a cosine of 0.84. Decoys at five times the
traffic leave the attack far above what popularity alone gives, and
\nudgeplain{}'s own differential privacy strengthens the attack while removing
most of the model's quality at our scale, for a reason that the ratio
$\users/\sqrt{\items}$ explains.

Third, composing creates a surface neither paper had. Once delivery feeds
training, a client's query is also a vote, and an unchecked vote of weight $10^6$
is available to anyone who calls the key generator directly. One round closes
inflation. Redistribution and the denial of service that full-domain evaluation
invites remain open, and both point to the same missing piece, an audit that a
key encodes a genuine point function.

The most useful practice in the project was measuring quality alongside cost.
A benchmark that counted only rounds and bytes would have reported that
revealing nothing costs 24\% of recommendation quality, when the real cost was a
silently clamping comparison range, and the true figure is two thousandths.

\paragraph{Next steps.} In order of value, the first is to batch the $\topk$ PIR
queries into one round trip, which our model says reverses the wide-area result.
The second is to drive the harvest from the networked servers, so that the loop
runs as a deployment and not only in tests. The third is the key audit, and the
fourth is private training at MovieLens-1M on a machine with the memory for it.
